% GENERATED FILE. Edit canonical lessons, then run npm run generate:books.
% canonical-source-hash: fnv1a64:284c40192e9a9a08
% canonical-lessons: GU-C452-khijavu, GU-C452-krodh, GU-C452-mijaj, GU-C452-ukalvu-seethe, GU-C452-kalah

\chapter{A Flash of Temper}
\label{ch:gu-a-flash-of-temper}
\hlchaptermodality{452}

\begin{chapteropening}
\textbf{By the end of this chapter:} \emph{I can say when someone loses their temper.}

Complete the last lesson of chapter 452: I can say when someone loses their temper.
\end{chapteropening}

\section[\texorpdfstring{\gu{ખીજાવું} (khījāvũ)}{ખીજાવું (khījāvũ)}]{\gu{ખીજાવું} (khījāvũ) — to get cross; to lose one's temper}

\label{lesson:GU-C452-khijavu}



\begin{quote}
Before the new one: say the Gujarati for love, then the Gujarati for to clasp.
\end{quote}

\subsection*{You'll want to know: \gu{ખીજાવું}}
\textbf{\gu{ખીજાવું}} — \emph{khījāvũ} — \textquotedblleft{}to get cross; to lose one's temper\textquotedblright{}.

\begin{quote}
\gu{દાદા} \gu{જલદી} \gu{ખીજાઈ} \gu{જાય} \gu{છે}.
\end{quote}

\begin{quote}
\emph{dādā jaldī khījāī jāy che.}
\end{quote}

\begin{quote}
Grandfather gets cross quickly.
\end{quote}

Past (he …): \gu{તે} \gu{ખીજાઈ} \gu{ગયો}, \emph{te khījāī gayo}.

\begin{cousinweb}
\textbf{\gu{ખીજાવું}} (\emph{khījāvũ}) is the \emph{-āvũ} form of the verb that CDIAL lists as Gujarati \emph{khījvũ}, \textquotedblleft{}to take offence\textquotedblright{}, under Sanskrit \emph{kṣīyatē}, \textquotedblleft{}is worn away, is injured\textquotedblright{} (CDIAL 3695).
\end{cousinweb}

\begin{grammarlens}[title={What you've built}]
One more word for describing what you lived and felt.
\end{grammarlens}

\subsection*{Your turn}
\begin{itemize}
\raggedright
  \item \emph{Say it:} \emph{khījāvũ}
  \item \emph{Say it:} \emph{khījāvũ}, once more
  \item \emph{Say it:} \emph{dādā jaldī khījāī jāy che.} — Grandfather gets cross quickly.
  \item \emph{Recall:} say the Gujarati for love, then the Gujarati for to clasp, then say \emph{khījāvũ} again
\end{itemize}

\begin{tcolorbox}[breakable,colback=teal!4,colframe=teal!35!black,title={Before you move on}]
What does \gu{ખીજાવું} mean? (\textquotedblleft{}To get cross; to lose one's temper\textquotedblright{}.) Say it once more.
\end{tcolorbox}
\section[\texorpdfstring{\gu{ક્રોધ} (krodh)}{ક્રોધ (krodh)}]{\gu{ક્રોધ} (krodh) — wrath; fierce anger}

\label{lesson:GU-C452-krodh}



\begin{quote}
Before the new one: say the Gujarati for to clasp, then the Gujarati for to get cross.
\end{quote}

\subsection*{You'll want to know: \gu{ક્રોધ}}
\textbf{\gu{ક્રોધ}} — \emph{krodh} — \textquotedblleft{}wrath; fierce anger\textquotedblright{}.

\begin{quote}
\gu{ક્રોધમાં} \gu{તેણે} \gu{ખોટો} \gu{જવાબ} \gu{આપ્યો}.
\end{quote}

\begin{quote}
\emph{krodhmā̃ teṇe khoṭo javāb āpyo.}
\end{quote}

\begin{quote}
In his anger he gave a false answer.
\end{quote}

A masculine noun.

\begin{cousinweb}
\textbf{\gu{ક્રોધ}} (\emph{krodh}) is a tatsama, Sanskrit \emph{krodha}, \textquotedblleft{}anger, wrath\textquotedblright{}, as Monier-Williams records it; everyday Gujarati says \textbf{\gu{ગુસ્સો}} (\emph{gusso}).
\end{cousinweb}

\begin{grammarlens}[title={What you've built}]
One more word for describing what you lived and felt.
\end{grammarlens}

\subsection*{Your turn}
\begin{itemize}
\raggedright
  \item \emph{Say it:} \emph{krodh}
  \item \emph{Say it:} \emph{krodh}, once more
  \item \emph{Say it:} \emph{krodhmā̃ teṇe khoṭo javāb āpyo.} — In his anger he gave a false answer.
  \item \emph{Recall:} say the Gujarati for to clasp, then the Gujarati for to get cross, then say \emph{krodh} again
\end{itemize}

\begin{tcolorbox}[breakable,colback=teal!4,colframe=teal!35!black,title={Before you move on}]
What does \gu{ક્રોધ} mean? (\textquotedblleft{}Wrath; fierce anger\textquotedblright{}.) Say it once more.
\end{tcolorbox}
\section[\texorpdfstring{\gu{મિજાજ} (mijāj)}{મિજાજ (mijāj)}]{\gu{મિજાજ} (mijāj) — temper; mood}

\label{lesson:GU-C452-mijaj}



\begin{quote}
Before the new one: say the Gujarati for to get cross, then the Gujarati for wrath.
\end{quote}

\subsection*{You'll want to know: \gu{મિજાજ}}
\textbf{\gu{મિજાજ}} — \emph{mijāj} — \textquotedblleft{}temper; mood\textquotedblright{}.

\begin{quote}
\gu{આજે} \gu{સાહેબનો} \gu{મિજાજ} \gu{ખરાબ} \gu{છે}.
\end{quote}

\begin{quote}
\emph{āje sāhebno mijāj kharāb che.}
\end{quote}

\begin{quote}
The boss is in a bad mood today.
\end{quote}

A masculine noun.

\begin{cousinweb}
\textbf{\gu{મિજાજ}} (\emph{mijāj}) is usually traced to Arabic \emph{mizāj}, \textquotedblleft{}temperament, mixture\textquotedblright{}, which reached Gujarati through Persian.
\end{cousinweb}

\begin{grammarlens}[title={What you've built}]
One more word for describing what you lived and felt.
\end{grammarlens}

\subsection*{Your turn}
\begin{itemize}
\raggedright
  \item \emph{Say it:} \emph{mijāj}
  \item \emph{Say it:} \emph{mijāj}, once more
  \item \emph{Say it:} \emph{āje sāhebno mijāj kharāb che.} — The boss is in a bad mood today.
  \item \emph{Recall:} say the Gujarati for to get cross, then the Gujarati for wrath, then say \emph{mijāj} again
\end{itemize}

\begin{tcolorbox}[breakable,colback=teal!4,colframe=teal!35!black,title={Before you move on}]
What does \gu{મિજાજ} mean? (\textquotedblleft{}Temper; mood\textquotedblright{}.) Say it once more.
\end{tcolorbox}
\section[\texorpdfstring{\gu{ઉકળવું} (ukaḷvũ)}{ઉકળવું (ukaḷvũ)}]{\gu{ઉકળવું} (ukaḷvũ) — to seethe; to boil}

\label{lesson:GU-C452-ukalvu-seethe}



\begin{quote}
Before the new one: say the Gujarati for wrath, then the Gujarati for temper.
\end{quote}

\subsection*{You'll want to know: \gu{ઉકળવું}}
\textbf{\gu{ઉકળવું}} — \emph{ukaḷvũ} — \textquotedblleft{}to seethe; to boil\textquotedblright{}.

\begin{quote}
\gu{તેનો} \gu{જવાબ} \gu{સાંભળીને} \gu{અરુણ} \gu{ગુસ્સાથી} \gu{ઉકળી} \gu{ગયો}.
\end{quote}

\begin{quote}
\emph{teno javāb sā̃bhḷīne Aruṇ gussāthī ukḷī gayo.}
\end{quote}

\begin{quote}
Hearing his answer, Arun seethed with anger.
\end{quote}

Past (it …): \gu{તે} \gu{પાણી} \gu{ઉકળી} \gu{ગયું}, \emph{te pāṇī ukḷī gayũ}.

\begin{cousinweb}
\textbf{\gu{ઉકળવું}} (\emph{ukaḷvũ}) goes back to a reconstructed \emph{*utkalati}, \textquotedblleft{}goes out or up\textquotedblright{}; CDIAL lists the Gujarati verb, \textquotedblleft{}to boil\textquotedblright{}, beside \textbf{\gu{ઉકાળવું}} (\emph{ukāḷvũ}), \textquotedblleft{}to cause to boil\textquotedblright{} (CDIAL 1716).
\end{cousinweb}

\begin{grammarlens}[title={What you've built}]
One more word for describing what you lived and felt.
\end{grammarlens}

\subsection*{Your turn}
\begin{itemize}
\raggedright
  \item \emph{Say it:} \emph{ukaḷvũ}
  \item \emph{Say it:} \emph{ukaḷvũ}, once more
  \item \emph{Say it:} \emph{teno javāb sā̃bhḷīne Aruṇ gussāthī ukḷī gayo.} — Hearing his answer, Arun seethed with anger.
  \item \emph{Recall:} say the Gujarati for wrath, then the Gujarati for temper, then say \emph{ukaḷvũ} again
\end{itemize}

\begin{tcolorbox}[breakable,colback=teal!4,colframe=teal!35!black,title={Before you move on}]
What does \gu{ઉકળવું} mean? (\textquotedblleft{}To seethe; to boil\textquotedblright{}.) Say it once more.
\end{tcolorbox}
\section[\texorpdfstring{\gu{કલહ} (kalah)}{કલહ (kalah)}]{\gu{કલહ} (kalah) — strife; a family quarrel}

\label{lesson:GU-C452-kalah}



\begin{quote}
Before the new one: say the Gujarati for temper, then the Gujarati for to seethe.
\end{quote}

\subsection*{You'll want to know: \gu{કલહ}}
\textbf{\gu{કલહ}} — \emph{kalah} — \textquotedblleft{}strife; a family quarrel\textquotedblright{}.

\begin{quote}
\gu{તે} \gu{ઘરમાં} \gu{રોજ} \gu{કલહ} \gu{થાય} \gu{છે}.
\end{quote}

\begin{quote}
\emph{te gharmā̃ roj kalah thāy che.}
\end{quote}

\begin{quote}
There is strife in that house every day.
\end{quote}

A masculine noun.

\begin{cousinweb}
\textbf{\gu{કલહ}} (\emph{kalah}) is a tatsama, Sanskrit \emph{kalaha}, \textquotedblleft{}strife, contention, quarrel\textquotedblright{}, as Monier-Williams records it.
\end{cousinweb}

\begin{grammarlens}[title={What you've built}]
One more word for describing what you lived and felt.
\end{grammarlens}

\subsection*{Your turn}
\begin{itemize}
\raggedright
  \item \emph{Say it:} \emph{kalah}
  \item \emph{Say it:} \emph{kalah}, once more
  \item \emph{Say it:} \emph{te gharmā̃ roj kalah thāy che.} — There is strife in that house every day.
  \item \emph{Recall:} say the Gujarati for temper, then the Gujarati for to seethe, then say \emph{kalah} again
\end{itemize}

\begin{tcolorbox}[breakable,colback=teal!4,colframe=teal!35!black,title={Before you move on}]
What does \gu{કલહ} mean? (\textquotedblleft{}Strife; a family quarrel\textquotedblright{}.) Say it once more.
\end{tcolorbox}
